\subsection{拓扑群的子群}

设 G 是一个拓扑群，H 是 G 的一个子群。由 (GT')，G 的拓扑在 H 上诱导的拓扑与 H 的群结构相容。这样在 H 上定义的拓扑群结构称为由 G 的结构诱导的结构。每当我们把 G 的子群 H 当作拓扑群考虑时，除另有明确说明外，所指的总是这个诱导结构。

命题 1 拓扑群 G 的子群 H 的闭包 \(\overline{H}\) 是 G 的一个子群。若 H 是 G 的正规子群，则 \(\overline{H}\) 也是。

若 \(a, b \in \overline{\mathbf{H}}\)，则 \(ab^{-1} \in \overline{\mathbf{H}}\)，因为映射 \((x, y) \to xy^{-1}\) 在 \(\mathbf{G} \times \mathbf{G}\) 上连续且把 \(\mathbf{H} \times \mathbf{H}\) 变换为 \(\mathbf{H}\)（第一章 § 2 第 1 目定理 1）。同样，映射 \(x \to axa^{-1}\) 的连续性表明：若 \(\mathbf{H}\) 是正规子群，则 \(\overline{\mathbf{H}}\) 是正规子群。

特别地，仅由 G 的单位元组成的集合 \(\{e\}\) 的闭包 N 是 G 的一个正规子群；且 \(N=\{e\}\) 当且仅当 G 是 Hausdorff 群（§ 1 第 2 目命题 2）。

命题 2 若 G 是 Hausdorff 拓扑群，则 G 的交换子群的闭包是 G 的交换子群。

由于命题 1，我们可以局限于 H 在 G 中稠密的情形。连续函数 xy 与 yx 在 H × H 上相等，因而由恒等延拓原理（第一章 § 8 第 1 目命题 2 推论 1），它们在 G × G 上相等。

命题 3 设 G 是 Hausdorff 拓扑群，M 是 G 的任一子集。则 G 中与 M 的每个元素交换的元素组成的集 M' 是 G 的一个闭子群。特别地，G 的中心在 G 中是闭的。

因为 \(\mathbf{M}'\) 是诸集合 \(\mathbf{F}_m(m \in \mathbf{M})\) 之交，其中 \(\mathbf{F}_m\) 是所有这样的 \(x \in \mathbf{G}\) 组成的集，它们使 \(xm = mx\)；而 \(\mathbf{F}_m\) 是闭的（第一章 § 8 第 1 目命题 2）。

命题 4 设 G 是一个拓扑群，H 是 G 的一个子群，它在 H 的某一点处是局部闭的（第一章 § 3 第 3 目定义 2）。则 H 在 G 中是闭的。

由平移，H 在其每一点处都是局部闭的，即 H 在 G 中是局部闭的。设 V 是 e 在 G 中的一个对称开邻域，使 \(V \cap H\) 在 V 中是闭的。若 \(x \in \overline{H}\)，则 xV 与 H 相交；若 \(y \in xV \cap H\)，则有 \(x \in yV\)，且 \(y(V \cap H) = (yV) \cap H\) 在 yV 中是闭的。但 x 属于 \((yV) \cap H\) 的闭包，故 \(x \in H\)。

推论 拓扑群的子群是开的，当且仅当它有一个内点。每个开子群都是闭的。

若子群 H 有一个内点，则由平移，H 的所有点都是内点，因而 H 是开的。推论的第二个断言是命题 4 的一个特殊情形。

命题 5 拓扑群 G 的子群 H 是离散的，当且仅当 H 有一个孤立点。Hausdorff 群的每个离散子群都是闭的。

若 H 是离散的，则 H 的每一点都是孤立点。反之，若 H 有一个孤立点，则由平移，H 的每一点都是孤立点，因此 H 是离散的。若 H 离散且 G 是 Hausdorff 群，则存在 e 的一个邻域 V 使 \(V \cap H = \{e\}\)；\(\{e\}\) 在 G 中是闭的，因而在 V 中更加是闭的，于是 H 在 e 处是局部闭的。故由命题 4，H 在 G 中是闭的。

注。设 H 是拓扑群 G 的任一子群。对每个 \(x \in \overline{H}\)，有 \(\overline{xH} = x\) . \(\overline{H} = \overline{H}\)，因为 x 平移是 G 到 G 上的同胚。换言之，对每个 \(x \in \overline{H}\)，xH 在 \(\overline{H}\) 中稠密。由此可得，若 H 不是闭的，则 \(\overline{H} \cap \mathrm{C}H\) 在 \(\overline{H}\) 中稠密。

